% ==================== 第 4 章 质量档案与迭代 ====================
\section{质量档案与迭代（人工视角）}

\noindent 分工说明：第 3 章 3.4 讲系统运行时\textbf{自动}发现与修正的机制；本章只收录\textbf{人工主动检查/测试发现的质量问题及其修复}，两章内容不混编。

\subsection{质量检查标准（团队口径声明）}

本节的“标准”是团队判断输出数据质量好坏的判据与口径，与 3.4.1 所述系统自动评估的\textbf{运行机制}互补：自动评估按来源逐条打分是“机器怎么量”，本节声明“量什么、什么叫对、谁来核”。两份口径全文同源、互相指路，不另起炉灶。

\subsubsection{检查维度与理由（评哪些方面、为什么）}

六个维度，直接对应赛道“查找完备、来源可溯、清洗可靠、便于分析”四评价点与多源整合加分项：

\begin{longtable}{@{}p{1.9cm}p{6.5cm}p{4.0cm}p{2.7cm}@{}}
\toprule
维度 & 检查内容 & 对应的赛道要求/加分项 & 与 3.4.1 自动评估键的对应 \\
\midrule
① 值正确性 & 数值与原文摘录逐字一致、字段归属正确（不错标字段）、不越物理合理范围；图表来源还要核坐标轴/图例/单位与读数是否相符 & 加分项“图表坐标轴或图例解析错误”；清洗可靠 & extraction\_quality、字段/单位校验 \\
② 完整性 & 预期应查到的性质与来源是否都覆盖；查询某性质却 0 条时是否显式呈现而非静默消失 & 加分项“缺失数据”；查找完备 & completeness \\
③ 一致性 & 单位写法、数值格式、重复记录是否被统一/去重；跨来源同量能否直接比较 & 加分项“重复数据、单位不一致”；清洗可靠 & consistency、format \\
④ 冲突处理 & 同一物理量的多来源分歧是否被识别、归因并保留标注，而非静默合并或丢弃 & 清洗可靠、来源可溯 & conflict\_risk、方差分析 \\
⑤ 来源可追溯 & 每条记录能否回到原始出处（论文页码+定位框+摘录，或星表行列位置），模型补充与规则换算能否与原始数据区分 & 来源是否可追溯 & source\_reliability、定位校验 \\
⑥ 输出可用性 & 数值/误差/单位是否机读、CSV/JSON 能否直接导入下游分析、前端与导出是否同源 & 便于后续分析 & 导出校验（quarantine 机制） \\
\bottomrule
\end{longtable}

\subsubsection{检查对象、样本与主体}

检查对象为当次查询产出的全部结构化记录；\textbf{评价主体分两层}——① 系统运行时对全量记录逐条自动评估（规则+模型，见 3.4.1），② 团队在报告与迭代期对输出做人工独立复核，抽样样本集中在三类高风险记录上：跨来源冲突组、标注“待核对”的记录、跨单位/跨写法记录（对照原文摘录与定位框逐条核验，见 4.2 记录与 4.3 迭代）。

\subsubsection{五类状态定义（正确 / 错误 / 缺失 / 冲突 / 无法判断）}

\begin{itemize}[leftmargin=1.2em,itemsep=2pt]
\item \textbf{正确}：数值与单位同原文摘录逐字一致，字段归属正确、落在合理范围，来源信息完整可回溯；
\item \textbf{错误}：与原文不符或字段归属错（错标字段、混入他对象值）、格式非法/类型混用、超出物理合理范围、图表读数与坐标轴/图例不符；
\item \textbf{缺失}：某来源预期应有而实际没有、或查询性质零产出——保留显式缺失标记并在质量报告中列出，不填充编造值；
\item \textbf{冲突}：同一物理量多来源取值明显分群（方法/时代/口径差异）——全部保留、逐条标注归因，不自动合并；
\item \textbf{无法判断}：实体归属存疑、异常无法验证、需要领域知识裁决——标注“待核对”或送人工，系统不擅断。
\end{itemize}

\subsubsection{口径一致性与留痕（修正前后同口径）}

机器修正（单位换算、格式规整、去重）与人工裁决（保留哪条、是否剔除）严格分开记录——每次修改写入前后对照（原值→新值、执行环节、原因），修正前后数据均可恢复；同一套五类定义贯穿自动评估、人工复核与 V1→V2 版本对比，不因修正放宽口径。此即模板 P5“修正前后同一口径”与 P12“保留修正前后真实差异”两问的合并回答。

\bigskip
\noindent 4.2 记录表与 4.3.2 第一版检查即按上述六维度与五类定义执行；系统自动评估的运行细节与综合分构成见 3.4.1。

\subsection{质量检查记录}

下表收录团队在\textbf{开发与评测期人工检查/主动测试发现}的质量问题及修复（自动运行机制内的检测不在此表，见 3.4）；六行按 4.1 六维度组织，问题按 4.1 五类状态归类；素材可回溯至代码注释事故档案与 git 提交记录（2026-07 至 09），V1→V2 案例级迭代已在 4.3 展示，此处不重复。

\begin{longtable}{@{}p{2.4cm}p{3.0cm}p{3.0cm}p{3.4cm}p{3.2cm}@{}}
\caption{人工检查发现的质量问题与修复}\label{tab:quality-records}\\
\toprule
实际质量问题（4.1 维度） & 系统如何发现 & 原因分析 & 实际修改的环节 & 修正后的结果 \\
\midrule
\endfirsthead
\toprule
实际质量问题（4.1 维度） & 系统如何发现 & 原因分析 & 实际修改的环节 & 修正后的结果 \\
\midrule
\endhead
\textbf{①③ 值正确性/一致性｜LLM 换算损坏真实数值}（实跑出现 0.118→0.0001、“100\_myr”→1e-7→0 式损坏，属\textbf{错误}类） & M45 真实查询输出与原文逐条比对时人工发现（2026-08-24） & 提取提示词未禁止换算且允许值内嵌单位，VLM 自行换算改写数值 & 三层同修：提取提示词硬规则（值与单位分离、禁止换算、逐字对应原文）→ 清洗规划器禁止生成换算代码（换算统一收敛到 unit\_converter 规则）→ 执行期可行域护栏（越界修改整工具回滚） & 同类损坏归零——V2 同查询 69 条值级修改全部为规整/格式清理，无换算改写（本行为 4.3.3“内容解析”行的成因纵深） \\
\textbf{⑤ 来源可回溯｜论文定位框不可用}（\textbf{错误}类） & 定位质量量化评测：旧整页 VLM 定位实测 IoU 中位数 0.00、$\geq$0.5 命中仅 2\%、59\% 中心偏差超 3 倍真值宽（2026-08-24） & 整页低分辨率下小数值漂移，且直接采信模型 bbox 无本地校验 & 重写混合双轨定位器：正文文本层矢量定位为主 + 表格整表宏定位 + 段落级兜底（零丢弃、零额外模型调用） & 正文定位命中率 97\%（矢量）、表格 100\%（整表），论文记录全部可点击回溯原文（V1 无此能力，见 4.3.2 ⑤） \\
\textbf{③ 一致性｜无实体锚点来源整表污染}（\textbf{错误}类） & 星表逐表人工审计（修 25 张表：剔除 key\_column 失效的 ocl/bax、补 I 前缀、去空格）+ 补充材料机制复盘 & 检索侧对无实体列（entity\_column 为空）的表格缺少保守策略，疑似扫描表/多天体表被整体纳入提取 & 检索：无 entity\_column 时仅 $\leq$3 行小表保留、否则整表丢弃；提取：VLM 提示词负向黑名单 + result\_builder 白名单强制校验 & 污染源被阻断——M45 案例 29 篇论文附带数据表经判定全部整表丢弃、零混入（见 3.3.2 补充材料行） \\
\textbf{④ 冲突处理｜人工裁决链路失效}（\textbf{无法判断}类出口不可用） & SN 2011fe 真实查询人工试用：14 项冲突堆叠成 17 张卡、逐条提问无限循环、取消后整批重审；2026-09-02 专项评审 & Source A/B 值未从 variance 回填（恒“?”）；前端 gate 误拦 human\_review 的 fields；逐条 interrupt 无批量协议；取消路径不清 pending；裁决理由链未贯通导出 & 人工审核批量重构：单次 interrupt 承载全批、逐项 1--5 强制裁决+可选理由、纯质量项仅 4/5、取消即清 pending、无效输入 $\leq$3 次全 skip；理由经 actions→source\_plan→data\_trace 贯通 & 同场景由约 3N+ 次 interrupt 降为固定 2 次；裁决理由随记录进入最终导出溯源（前端血缘可见）；无再发堆叠（见 2.4 人机协同承诺的落地） \\
\textbf{②⑥ 完整性/输出可用性｜Web 交付链路断链族}（\textbf{错误}类静默失效） & 2026-08-12 Web 专项审计：8 路 finder + 8 路对抗验证 + 动态探针，确认 60 条真实问题（3 critical） & 前后端契约字段不一致（task.id vs task\_id、source\_type 读错键）、summary 读错字段层级、图证 URL 路径缺段、导出落随机目录与任务解耦、SSE 断线全量重放致消息/澄清重复 & 前端与后端按契约同源修复，分 14 批处理（P0→P15 含回放验收） & 审计 60 条全部闭环：任务总结与数据同源、图证可加载、导出与任务绑定、断线续播不重弹——3 个 critical 级静默错误数据路径清除 \\
\textbf{① 值正确性｜人工出口被静默跳过}（\textbf{无法判断}类出口缺失） & SN 2011fe 实测 + 代码评审：checkpointer 缺失时质量子图 interrupt 走降级分支，人工审核被整条跳过且无提示 & web 主图 checkpointer 配置链取不到（config 为空）→ HITL 静默降级为“无人工”路径 & 进程级全局共享 checkpointer 兜底（两处同源修复）；纯质量类待审项（QHR）的降级策略显式化（skip=原样交付） & 需裁决时必弹人工卡、不再静默跳过；确属纯质量类则按显式规则降级，数据原样交付且留痕 \\
\bottomrule
\end{longtable}

\noindent 模板 P12 表外四问（哪些问题触发重新查找来源 / 哪些触发重新解析、字段对齐或单位换算 / 哪些必须交研究者或人工审核 / 如何保留修正前后真实差异）本质描述的是\textbf{流程内自动评价的路由机制}，已由 3.4 质量管线相应小节覆盖（重查来源 → 3.3.1 缺失处理与收尾；重解析/对齐换算复验 → 3.4.1/3.4.2；交人工 → 3.4.3 存疑处置与 HR；留痕 → 3.4.2 修改记录与 3.4.4 输出），第 4 章不再重复。

\subsection{版本迭代演示：V1 → V2}

V1 = 系统早期版本 2026-08-21 对 M45 同查询的真实输出（任务 906f3c93，165 条记录/30 来源，见 4.3.1）；V2 = 定稿版 2026-08-28 的同查询真实运行（任务 1e6c3e2e，103 条记录/30 来源，即样例库“08-28 版”）。本报告将模板 P16（第一版输出与质量检查）与 P17（第二版输出与迭代变化）合成一节，展示“按 4.1 标准检查 V1 → 逐项修改 → V2 复验”的一版完整迭代。

\subsubsection{第一版实际输出（V1）}

V1 = 系统早期版本对 M45 同查询的真实输出：\textbf{165 条记录 / 30 个来源}（MWSC 星团目录 + 29 篇论文）/ 5 项性质（distance 59、dist\_modulus 30、cluster\_age 29、fe\_h 28、parallax 19）；提取方式：论文正文 107 条、论文表格 54 条、数据库 4 条。下表为\textbf{原始输出直引}的代表性记录（值、单位、出处均未修饰）——它们就是 4.3.2 检查的原始对象：

\begin{longtable}{@{}p{3.3cm}p{4.2cm}p{1.7cm}p{3.9cm}@{}}
\toprule
对象 / 出处（真实定位） & 记录内容（直引原文记录） & 来源 & 人工初看印象（按 4.1 状态预判） \\
\midrule
MWSC 星团目录（J/A+A/558/A53，表 catalog、行 0305） & cluster\_age = 8.1499…（原始列 logt）log(yr)；distance = 130 pc 等 4 条目录记录，原始列名 logt/d/MOD/[Fe/H] 随记录保留 & 数据库查询 & 对数写法单位与论文体系异类 → 待检查 \\
1999ApJ…523..328N 第 13 页正文（bbox 可定位） & distance = “130.7 $\pm$ 11.1 pc” & 正文提取 & 值内嵌单位（字段另标 pc）→ 疑似脏格式 \\
2006A\&A…446..611F 第 9 页表格 & cluster\_age = “70-100”，单位 $10^{6}$ yr & 表格提取 & 区间写法 + 第三套单位 → 疑似脏格式 \\
2006A\&A…446..611F 第 9 页表格（另一行） & cluster\_age = “70$\times10^{6}$ - 100$\times10^{6}$”（单位 yr） & 表格提取 & 值与单位双写、范围写法 → 疑似脏格式 \\
2018ApJ…863..179S 表格 & cluster\_age = “120 Myr”（字段单位同为 Myr） & 表格提取 & 值内嵌单位且与单位列重复 → 疑似脏格式 \\
2018ApJ…863...67G 第 3 页表格 & dist\_modulus = “0.1054” mag（同表另一行 5.64 mag） & 表格提取 & 与同表 5.64 mag 并列、量级悬殊 → 疑似错标字段（错误） \\
2018ApJ…863...67G 表格 & cluster\_age = “$112^{+2}_{-26}$” Myr & 表格提取 & 不对称误差 LaTeX 写法 → 需规整 \\
2004A\&A…418L..31M & distance = “116 $\pm$ 3.3 pc”（同文还报 129$\pm$3、135$\pm$2 pc 等不同口径） & 正文提取 & 一源多口径并存 → 待核对 \\
\bottomrule
\end{longtable}

\subsubsection{第一版质量检查：按 4.1 标准对 V1 实施的第一轮检查}

检查对象 = V1 全部 165 条记录（全量，不抽样）；判定依据 = 4.1 六维度与五类定义；检查方式 = 系统自动评估（评估发现问题 102 项、方差组 67 项、异常 18 项）+ 人工对照原文摘录与同表邻值复核。结果即 4.3.3 修改表的输入，前后口径一致。

\begin{longtable}{@{}p{1.9cm}p{3.1cm}p{6.4cm}p{3.4cm}@{}}
\toprule
4.1 维度 & 检查了哪些结果、如何判断 & 第一版检查结果（真实数据） & 问题归类与去向 \\
\midrule
① 值正确性 & 全量 165 条逐条核字段归属与合理范围（对照原文摘录、同表邻值） & 发现问题 102 项；代表例：2018ApJ…863...67G 表中 dist\_modulus 0.1054 mag 与同表 5.64 mag（真距离模数）并列——0.1054 实为该表其它列的消光量，字段标错 & \textbf{错误}（错标字段）；越界值当时无护栏拦截 → 需 V2 加执行期拦截 \\
② 完整性 & 按性质清单检查 5 项性质是否都有产出 & 5/5 性质均有记录（min parallax 19 条），未发现零产出；无缺失需处理 & 未见异常（如实记录） \\
③ 一致性（单位/格式/重复） & 数值格式、单位写法、重复条目 & cluster\_age 单位 \textbf{4 种写法并存}（Myr 19 / yr 8 / log(yr) 1 / $10^{6}$ yr 1）；值级脏格式十余条——值内嵌单位（“120 Myr”、“125 $\pm$ 8 Myr”、“~ 100 Myr”）、范围/区间写法（“70-100”、“70$\times10^{6}$-100$\times10^{6}$”、“~ 120-130 Myr”）、不对称误差（“$112^{+2}_{-26}$”、“132 +26 $-$27”）、科学计数（“1.2e8”）；同源近似重复（“112”$\times$3、“118$\pm$6”三种写法$\times$3）；distance 的 $\pm$ 空格与写法亦不齐（“134$\pm$3” vs “136.2$\pm$1.2”） & \textbf{错误}（脏格式不可直接比较）→ 需 V2 数值与单位分离硬约束 + 去重 \\
④ 冲突与异常 & 方差分析输出 & 已检出跨源方差组 67 项、异常 18 项（2 个来源被路由至冲突路径），但 V1 阶段只有初级的“冲突解决 6 项”式处理，\textbf{无系统归因、无保留标注}——真分歧与“单位写法不同造成的伪分歧”混在一起 & \textbf{冲突}/无法判断混杂 → 需 V2 系统化归因 \\
⑤ 来源可回溯 & 每条记录是否带来源与定位 & 165 条全部带来源与定位（数据库行列 / 论文页码+定位框），形式上可点击回溯 & 定位框由模型直接给出、未经本地校验，精度存疑 → 需 V2 本地逆向定位复核 \\
⑥ 输出可用性 & 是否可直接导入下游分析 & 脏格式与单位不统一使记录无法直接计算/比较；前端展示的是清洗前数据，与导出文件不同源 & \textbf{不可用} → 需 V2 统一数据源 \\
\bottomrule
\end{longtable}

\begin{itemize}[leftmargin=1.2em,itemsep=2pt]
\item \textbf{第一版是否能够直接用于后续科研分析，依据是什么：}不能。① 值级脏格式（值内嵌单位、范围写法、不对称误差）无法直接计算——“120-130 Myr”、“70$\times10^{6}$-100$\times10^{6}$”之类字符串不能进统计程序；② 距离模数含错标字段值（0.1054）；③ 单位 4 种写法并存无法跨来源比较。直接用于分析前需一条条人工清洗换算。
\item \textbf{第一版需要修改的点（即 4.3.3 各行的修改输入）：}① 提取阶段缺少数值与单位分离硬约束 → 值内嵌单位/范围写法；② 年龄单位未按统一规则换算；③ 错标字段与超范围值缺乏执行期拦截；④ 原文定位依赖模型定位框、精度不足；⑤ 前端展示与导出数据不同源。
\end{itemize}

\subsubsection{第一版与第二版的实际变化（V1 vs V2）}

V2 为系统定稿后（2026-08-28，任务 1e6c3e2e）同查询真实运行，共 103 条记录、30 个来源、78 张图证、综合质量分 0.85——定稿版提取约束与清洗规则增强，记录数较 V1（165 条）收敛，去重与低置信过滤更严格（4.3.2 ⑥ 类问题亦在此版修复）。下列各行第一版结果均对应 4.3.1/4.3.2 的真实输出与检查，无变化处如实注明“无”。

\begin{longtable}{@{}p{1.8cm}p{3.6cm}p{2.4cm}p{3.5cm}p{3.2cm}@{}}
\toprule
变化内容 & 第一版结果 & 发现的问题 & 实际修改的方法 & 第二版结果 \\
\midrule
\endfirsthead
\toprule
变化内容 & 第一版结果 & 发现的问题 & 实际修改的方法 & 第二版结果 \\
\midrule
\endhead
数据源或查找范围 & 30 个来源（29 篇论文 + MWSC 星团目录） & 无 & 无 & 无变化（同为 30 个来源） \\
内容解析 & cluster\_age 单位 4 种写法并存、值级脏格式十余条（“120 Myr”值内嵌单位、“70$\times10^{6}$ - 100$\times10^{6}$”、“$112^{+2}_{-26}$”等，见 4.3.1/4.3.2 ③）；原文定位由模型直接给出，精度差（见 4.3.2 ⑤） & 提取阶段缺少数值与单位分离硬约束；模型定位不可靠 & 提取提示词增加数值与单位分离、禁止换算硬规则；原文定位改为按页面文本层逆向定位，表格改为整表定位 & 脏格式经 69 条值级修改全部规整留痕（如 “134.4 +2.9 $-$2.8”→134.4、“7.48 $\pm$ 0.03”→7.48）；论文记录全部附带可靠原文定位（正文矢量定位、表格整表定位，零额外调用成本） \\
重复、缺失或冲突处理 & 67 组方差/18 项异常仅初级处理（冲突解决 6 项）、无归因标注（见 4.3.2 ④）；错标字段值（dist\_modulus 0.1054）未被拦截 & 超范围修改缺乏执行期拦截；冲突无系统归因 & 增加执行期越界拦截，超出合理范围的修改整体回滚；冲突处理增强（归因分类+保留标注） & 31 项跨源冲突归因（26 项自动、0 项需人工）、剩余问题 0；0 条越界值进入输出 \\
结构化输出 & 前端表格展示的是清洗前数据，与导出文件不一致 & 前端展示与导出数据不同源 & 前端记录接口改为返回清洗后的最终数据，与导出文件同源 & 前端展示与导出文件一致 \\
\bottomrule
\end{longtable}

\subsubsection{第二版结果与未解决问题}

\begin{itemize}[leftmargin=1.2em,itemsep=2pt]
\item \textbf{第二版相较第一版实际改善了什么：}数值与单位分离的脏格式经 69 条值级修改全部规整留痕；异常（26 项）与冲突（31 项）全部主动标记而非静默处理；错误标字段与越界修改被拦截（0 条越界值进入输出）；原文定位全部转为可靠的双轨定位；前端展示与导出数据一致。
\item \textbf{哪些问题没有解决或出现了新的代价：}跨来源单位差异收规则限制，依然无法保证全局统一；范围写法记录（如 “75-100 Myr”、“fe\_h -0.14 to +0.13”）按设计保留待核对。
\item \textbf{第二版数据能够支持什么后续分析：}支持星团年龄、距离、金属丰度的多源对比与参数汇总（现代测距值收敛区间 134.4--136.2 pc，年龄 100--130 Myr，金属丰度 $-$0.03~+0.03 dex），冲突标注支持研究者逐条裁决。
\item \textbf{第二版是否达到团队设定的质量要求，依据是什么：}达到。综合质量分 0.85，清洗与冲突处理循环 1 轮后剩余问题为 0（31 项冲突 26 项自动归因、0 项需人工）。
\end{itemize}

% ==================== 第 5 章 对照与消融 ====================
\section{对照与消融}

\subsection{团队实际完成的比较}

四组对照均真实执行（含人工整理对照组），保持同一查询（M45）、同一 5 项标准性质与同一评价口径。

\begin{longtable}{@{}p{2.3cm}p{1.9cm}p{2.4cm}p{3.0cm}p{2.9cm}p{2.2cm}@{}}
\caption{实际完成的对照与消融}\label{tab:ablation}\\
\toprule
实际比较对象 & 保持相同的数据需求 & 采用的评价方法 & 本作品结果 & 对照结果 & 结论与代价 \\
\midrule
\endfirsthead
\toprule
实际比较对象 & 保持相同的数据需求 & 采用的评价方法 & 本作品结果 & 对照结果 & 结论与代价 \\
\midrule
\endhead
对照组 1：人工整理 & 同一查询、同一 5 项性质 & ① 记录与来源条数 ② 提取时间长度 ③ 正确率 & ① 103 条 / 30 篇来源 ② 耗时 21 分钟 ③ 质量分 0.8502，有 2--3 处错误 & ① 34 条 / 12 篇来源 ② 耗时 3.5 小时 ③ 基本无错误 & 人工提取条数与来源显著少于本产品，时间显著多于本产品。在质量上几乎无差距；代价：专业从业人员提取的质量可能好于产品 \\
对照组 2：无质量闭环单轮提取（全链路消融：仅移除质检环） & 同上 & ① 值可计算性 ② 单位一致性 ③ 异常与冲突处理 ④ 可信度评估 & ① 103 条经 69 条值级修改，大部分全部规整留痕 ② 物理单位（如年龄）基本归一，仅存在少数未匹配规则的情况保留真实单位全量评估 ③ 26 项异常主动标记、31 项跨源冲突归因 ④ 质量分 0.8502 & ① 110 条不确定度/区间/双$\pm$ 文本混排（如 120-130 Myr、136.8 $\pm$ 6.4 pc）不可直接计算 ② age 三单位未统一数量显著多于完整版 ③ 跨源差异无识别、没有处理冲突数据 ④ 无质量分、无可信度标注 & 质检环有效解决了裸提取“格式混排不可算、单位混乱、冲突无归因”的问题，是数据达到科研可用门槛的核心模块。代价：质检增加约 6 分钟、64 次 LLM 调用（端到端 21 分钟） \\
对照组 3：单一来源（仅 MWSC 星团目录，无论文及补充材料源） & 同上 & ① 性质全面性 ② 数据多样性 ③ 跨源验证能力 & ① 覆盖度 5/5 性质 ② 数据 103 条 ③ 多源收敛：Gaia 135.15 / 恒星双胞胎法 134.8 / 移动星团法 134.4 pc & ① 性质缺 parallax ② 数据仅有 4 条 ③ 130 pc 单值无相互印证 & 无论文源下，无法满足每个性质的数据要求；数据量不足，无法交叉验证。代价：检索时间增长约 10 分钟 \\
对照组 4：通用检索（裸大模型一次问答，无检索/无提示词工程） & 同上 & ① 结构化程度 ② 来源可验证性 ③ 单位规范 & ① 每条值/单位分离 ② 逐值表格行列/bbox 框选定位 ③ JSON/CSV 直读入分析 & ① 1 段自然语言、3 个区间值 ② 零出处、零定位、无溯源可能 ③ 误差以自然语言混排 & 结构化+可验证替代自由文本。代价：端到端 21 分钟，不如直接问答快 \\
\bottomrule
\end{longtable}

\subsection{结果说明}

\begin{itemize}[leftmargin=1.2em,itemsep=2pt]
\item \textbf{科研数据适用性方面，本作品实际改善了什么：}解决了原始提取数据“文本混排、不可计算、无法直接做科研统计”的痛点。质检环前，66 项数据内嵌不确定度（如 134.8 $\pm$ 1.7 pc）、范围区间（如 100--120 Myr）及多重单位，下游分析脚本无法直接读取；质检后，数值与误差分离、物理单位基本统一为标准机读格式。M45 的三项核心参数得以直接收敛出统计结果（距离 134.4--136.2 pc、年龄 100--130 Myr、金属丰度 $-$0.03~+0.03 dex），并附带全链路质量分（0.8502），为下游科研绘图与科学分析提供了开箱即用、带置信度依据的标准数据集。
\item \textbf{数据处理方法方面，本作品实际改善了什么：}构建了全流程留痕的质检闭环，改变了裸提取缺乏校验手段的缺陷：① \textbf{规范化留痕}：完成 69 处数值级清洗与拆解（如 134.4 +2.9 $-$2.8 拆解为中心值 134.4 与独立误差字段），全程记录前后对照日志；② \textbf{离群主动识别}：通过统计模型检出 26 项异常（22 项确认异常、4 项物理边界、0 误报），成功将 MWSC 目录早期的 130 pc 识别为统计离群值；③ \textbf{跨源差异归因}：对 31 项跨源数值冲突自动判定原因，识别出由测量方法不同（如主序拟合法 vs Gaia 视差法）导致的系统偏差；④ \textbf{护栏阻断机制}：通过物理量纲校验，彻底拦截“消光值误识为距离模数”等跨语义错位，保证输出端 0 越界值。
\item \textbf{结果质量和复用方面，本作品实际改善了什么：}① \textbf{证据链闭环}：103 条数据中 99 条（96.1\%）精准绑定 PDF 物理页码与 Bounding Box 框选坐标，科研人员点击任意数值即可一键直达原文图文核对；② \textbf{数据资产标准化}：导出与看板同源的结构化 JSON/CSV，完整封装标准化数值、不确定度、物理单位及文献 DOI/BibTeX 等元数据，完全满足科研成果公开发布与代码二次复现的要求。
\item \textbf{没有改善的部分及增加的成本（如 Token 开销、耗时）：}① \textbf{计算与耗时代价}：质检闭环增加了 64 次 LLM 调用与约 6 分钟耗时（端到端总耗时约 21 分钟，质检开销占比约 24\%）；② \textbf{极端复杂区间的科学边界}：对部分学界本身存在争议的复杂范围写法（如金属丰度 $-$0.14 to +0.13 dex），系统选择保留原文并打上提示标签，将科学裁决权保留给学者，不做武断的截断处理；③ \textbf{上游依赖与检索漂移}：对于深层天体物理上下文（如双星系统的成员消歧），系统目前仍依赖上游解析，质检管线未做深层动力学推理。
\end{itemize}
